\section{Raviolo restriction: Čech/Thom--Sullivan model and coinvariants}
\label{sec:raviolo-expanded}

The \v{C}ech/Thom--Sullivan totalization of the raviolo factorization algebra (\S\ref{sec:raviolo}) computes derived global sections; the resulting homotopy coinvariant spaces furnish the generic raviolo coinvariant package, and on finite-type comparison surfaces their linear duals recover the corresponding conformal blocks.

\subsection{Raviolo geometry and derived global sections}
\label{subsec:raviolo-geometry}
Let $\mathrm{Rav}:=D_+\cup_{D^\times}D_-$ denote the raviolo obtained by gluing two copies of the formal disk along the punctured disk. Cover by $\{D_+,D_-,D^\times\}$ and write the Čech nerve.
\begin{definition}[Raviolo factorization algebra]
\label{def:raviolo-FA}
Given a $W(\mathsf{SC}^{\mathrm{ch,top}})$–algebra $(A_{\mathsf{ch}},A_{\mathsf{top}})$, define the raviolo factorization algebra $\mathsf V_{\mathrm{rav}}$ by restricting to time–slices as in Section~\ref{sec:raviolo} and glue along $D^\times$ using the $E_1$–homotopy. Then
\[
\Gamma(\mathrm{Rav},\mathsf V_{\mathrm{rav}})
:= \Tot\bigl(\check{C}^\bullet(\{D_+,D_-,D^\times\};\mathsf V_{\mathrm{rav}})\bigr)
\]
with Thom–Sullivan realization yields the derived global sections.
\end{definition}
\begin{proposition}[Čech exactness; \ClaimStatusProvedElsewhere]
\label{prop:Cech-exactness}
For holomorphic–topological $\mathsf V_{\mathrm{rav}}$, the Čech/Thom–Sullivan totalization computes derived global sections on $\mathrm{Rav}$.
\end{proposition}
\begin{proof}
Use cover descent and the standard argument for factorization algebras; cf.\ \cite{CG17,GW23}.
\end{proof}

\subsection{Coinvariants and conformal blocks}
\label{subsec:coinvariants}
\begin{definition}[Coinvariants]
\label{def:coinvariants}
For a raviolo VA $\mathsf V_{\mathrm{rav}}$, define the coinvariants on a test curve $X\to\mathrm{Rav}$ by the homotopy quotient
\[
\mathrm{Coinv}_X(\mathsf V_{\mathrm{rav}}) := \bigl(\Gamma(X,\mathsf V_{\mathrm{rav}})\bigr)_{G}^{\mathrm{h}}
\]
with respect to the global symmetries $G$ acting by currents (where $G$ is a Lie algebra acting on~$\mathsf V_{\mathrm{rav}}$ by chiral currents, when such an action is part of the data). Conformal blocks are the linear duals when finite type.
\end{definition}

\begin{remark}
The holomorphic–topological pairing ensures that distributions supported on $D^\times$ implement the usual vertex–algebra state–field correspondence after passing to cohomology.
\end{remark}

\subsection{Shifted Poisson vertex structure on cohomology}
\label{subsec:PVA-cohomology-expanded}
Let $H^\bullet(\mathsf V_{\mathrm{rav}})$ be the BV cohomology. By Section~\ref{sec:Ainfty-to-PVA} and Theorem~\ref{thm:raviolo-PVA}, we have:

\begin{theorem}[Cohomological PVA on raviolo; \ClaimStatusProvedElsewhere]
\label{thm:raviolo-PVA-expanded}
\textup{(Reproduced from Theorem~\ref{thm:raviolo-PVA} for the reader's convenience.)}
$H^\bullet(\mathsf V_{\mathrm{rav}})$ is a $(-1)$–shifted Poisson vertex algebra; its $\lambda$–bracket is given by the residue of the antisymmetric singular part of $m_2$, and its product is the symmetric regular part of $m_2$. Higher $m_{k\ge3}$ vanish in cohomology by Stokes/AOS identities.
\end{theorem}

\begin{corollary}[Agreement with HT constructions; \ClaimStatusConditional]
\label{cor:agreement-HT}
For HT twists of 3d $\mathcal N=2,4$ theories whose observables lie in the scope of Theorem~\ref{thm:physics-bridge}, $H^\bullet(\mathsf V_{\mathrm{rav}})$ agrees with the PVA extracted by the secondary products in \cite{OY20,CDG20} and the 2--shifted Poisson structures in \cite{GW23,GRW21}.  The two shifts are reconciled by the relationship between the derived algebraic geometry perspective (which counts the shift of the symplectic form on the derived stack, giving~$+2$) and the vertex algebra perspective (which counts the intrinsic degree of the $\lambda$-bracket as an operation on the underlying graded space, giving~$-1$).  The total shift $2 + (-1) = 1$ accounts for the passage from the stack-theoretic to the operadic description, with the~$+1$ absorbed by the parity of the raviolo.
\end{corollary}
